\documentclass[11pt,a4paper]{article}
% main (standalone preamble for editing)
% vim: nowrap: spell:
% ══════════════════════════════════════════════════════════════════════════════
% Speed Maths --- shared preamble
% Included by every sheet and answer file via:  % main (standalone preamble for editing)
% vim: nowrap: spell:
% ══════════════════════════════════════════════════════════════════════════════
% Speed Maths --- shared preamble
% Included by every sheet and answer file via:  % main (standalone preamble for editing)
% vim: nowrap: spell:
% ══════════════════════════════════════════════════════════════════════════════
% Speed Maths --- shared preamble
% Included by every sheet and answer file via:  \input{../../shared/preamble}
% Editing THIS file changes the look of every sheet/answer across every pillar.
% ══════════════════════════════════════════════════════════════════════════════

\usepackage[a4paper, top=25mm, bottom=25mm, left=22mm, right=22mm]{geometry}
\usepackage{amsmath, amssymb}
\usepackage{enumitem}
\usepackage{booktabs}
\usepackage{array}
\usepackage{parskip}
\usepackage{microtype}
\usepackage{fancyhdr}
\usepackage{titlesec}
\usepackage{mdframed}
\usepackage{xcolor}
\usepackage[hidelinks]{hyperref}
\usepackage{graphicx}
\usepackage{tikz}
\usepackage{pgfplots}
\pgfplotsset{compat=1.18}

% ── Colours ───────────────────────────────────────────────────────────────────
\definecolor{darkgray}{gray}{0.15}
\definecolor{medgray}{gray}{0.45}
\definecolor{lightgray}{gray}{0.93}
\definecolor{boxgray}{gray}{0.88}

% ── Section formatting ──────────────────────────────────────────────────────────
\titleformat{\section}[block]
  {\normalfont\large\bfseries}{}
  {0pt}{}[\vspace{2pt}\hrule\vspace{6pt}]
\titlespacing*{\section}{0pt}{14pt}{4pt}

% ── List formatting ─────────────────────────────────────────────────────────────
\setlist[enumerate,1]{label=\textbf{\arabic*.}, leftmargin=2em, itemsep=10pt}

% ── Branding constants (edit here, not per-file) ────────────────────────────────
\newcommand{\SpeedExamLine}{\textbf{Speed Maths:} Competition \& Exam Prep}
\newcommand{\SpeedCredit}{Series by \href{https://www.linkedin.com/in/cerealdev/}{David Akinyele-Aje}}

% ── Header / footer ──────────────────────────────────────────────────────────────
% Usage (once, after \input{../../shared/preamble} and before \begin{document}):
%   \SpeedHeader{Algebra}{3}
% First arg  = pillar name shown as "Speed <Pillar> --- Daily Drill"
% Second arg = worksheet number
\pagestyle{fancy}
\newcommand{\SpeedHeader}[2]{%
  \fancyhf{}%
  \renewcommand{\headrulewidth}{0.4pt}%
  \setlength{\headheight}{14pt}%
  \fancyhead[L]{\small\textbf{Speed #1 --- Daily Drill}}%
  \fancyhead[R]{\small Worksheet \##2}%
  \fancyfoot[L]{\small \SpeedExamLine}%
  \fancyfoot[C]{\small\thepage}%
  \fancyfoot[R]{\small \SpeedCredit}%
  \renewcommand{\footrulewidth}{0.4pt}%
}

% ── Title block (used at the top of every sheet AND answer file) ───────────────
% Usage: \SpeedTitleBlock{Daily Algebra Drill \#3}{Jane Doe}
% %% AGENTS: When generating a new sheet, ask the human contributor for the name they want credited in argument 2.
% For answer files, pass the "--- Answers \& Investigations" suffix in the arg.
\newcommand{\SpeedTitleBlock}[2]{%
  \begin{center}
    {\LARGE\bfseries #1}\\[4pt]
    {\normalsize\itshape Sheet by #2}\\[6pt]
    \hrule\vspace{4pt}
    \begin{tabular}{llcc}
      \toprule
      \textbf{Section} & \textbf{Topic} & \textbf{Questions} & \textbf{Target time}\\
      \midrule
      A & Rapid Recognition          & 10 & 2 min 30 sec \\
      B & Manipulation Drills        & 10 & 8 min        \\
      C & Substitution \& Structure  & 8  & 10 min       \\
      D & Challenge                  & 5  & 15 min       \\
      \bottomrule
    \end{tabular}
    \vspace{4pt}
    \hrule
  \end{center}
}

% ── Sheet metadata: topic + the day's new tools ────────────────────────────────
% Usage (immediately after \SpeedTitleBlock, sheets only):
%   \SpeedMeta{Expanding, factorising, and the identity toolkit}
%             {difference of squares, sum/product form, completing the square}
%
% Arg 1 = the sheet's topic in one line. The website lists this instead of
%         "Sheet 03", so write the thing a reader would actually search for.
% Arg 2 = comma-separated, at most five: the tools this sheet introduces that
%         earlier sheets in the pillar did not. These become the website's tags.
%         Leave out anything the reader already met — the tags are meant to say
%         what is new today, not to summarise the sheet.
%
% Both are parsed straight out of this file by tools/build_website.py, so the
% sheet stays the single source of truth and the site cannot drift from it.
%
% This renders NOTHING. It is a declaration for the build script, not a piece of
% the sheet's layout: adding it to a sheet must not change that sheet's PDF, so
% the 25 existing sheets can be annotated without a single one of them being
% reissued. If the toolkit should also appear on the page, that is a separate
% decision about the sheet design — make it deliberately, here, not by leaving
% this macro printing something nobody asked it to.
\newcommand{\SpeedMeta}[2]{}

% ── Answer / Method / Investigate-further boxes (answer files only) ────────────
\newmdenv[
  backgroundcolor=lightgray,
  linecolor=medgray,
  linewidth=0.5pt,
  leftmargin=0pt, rightmargin=0pt,
  innerleftmargin=8pt, innerrightmargin=8pt,
  innertopmargin=5pt, innerbottommargin=5pt,
  skipabove=4pt, skipbelow=4pt
]{ansbox}

\newmdenv[
  backgroundcolor=white,
  linecolor=darkgray,
  linewidth=0.8pt,
  leftline=true, rightline=false, topline=false, bottomline=false,
  leftmargin=0pt, rightmargin=0pt,
  innerleftmargin=8pt, innerrightmargin=6pt,
  innertopmargin=4pt, innerbottommargin=4pt,
  skipabove=4pt, skipbelow=6pt
]{invbox}

\newcommand{\ans}[1]{%
  \begin{ansbox}
    \textbf{Answer:} #1
  \end{ansbox}}

\newcommand{\method}[1]{%
  {\small\textbf{Method:} #1}\par}

\newcommand{\inv}[1]{%
  \begin{invbox}
    \textbf{\small Investigate further:} {\small #1}
  \end{invbox}}

% ── Misc helpers ────────────────────────────────────────────────────────────────
\newcommand{\qn}[1]{\textbf{#1.}\;}

% Mistake-linking: attach a Top 5 Patterns / Common Traps bullet to the question
% label(s) in this sheet where it actually shows up, e.g. \seealso{B6, D2}
\newcommand{\seealso}[1]{\hfill\textit{\footnotesize(see #1)}}

% ── Graph options macro ─────────────────────────────────────────────────────────
% Usage: \GraphOptions{tikz code for A}{tikz code for B}{tikz code for C}{tikz code for D}
\newcommand{\GraphOptions}[4]{%
  \begin{center}
    \begin{tabular}{cc}
      \textbf{A)} & \textbf{B)} \\
      #1 & #2 \\[10pt]
      \textbf{C)} & \textbf{D)} \\
      #3 & #4
    \end{tabular}
  \end{center}
}

% ── Closing motivational rule + cross-promo footer (sheets only) ────────────────
% Usage: \SpeedClosing{"Quote text."}
\newcommand{\SpeedClosing}[1]{%
  \vspace{20pt}
  \begin{center}
  \rule{0.6\textwidth}{0.4pt}\\[4pt]
  {\small\itshape #1}\\[4pt]
  \rule{0.6\textwidth}{0.4pt}
  \end{center}
  \begin{center}
  {\small Found a mistake or want to contribute a sheet?}\\
  {\small Visit the open-source project at \href{https://github.com/The-CerealDev/speed-maths}{github.com/The-CerealDev/speed-maths}}
  \end{center}
}

% Editing THIS file changes the look of every sheet/answer across every pillar.
% ══════════════════════════════════════════════════════════════════════════════

\usepackage[a4paper, top=25mm, bottom=25mm, left=22mm, right=22mm]{geometry}
\usepackage{amsmath, amssymb}
\usepackage{enumitem}
\usepackage{booktabs}
\usepackage{array}
\usepackage{parskip}
\usepackage{microtype}
\usepackage{fancyhdr}
\usepackage{titlesec}
\usepackage{mdframed}
\usepackage{xcolor}
\usepackage[hidelinks]{hyperref}
\usepackage{graphicx}
\usepackage{tikz}
\usepackage{pgfplots}
\pgfplotsset{compat=1.18}

% ── Colours ───────────────────────────────────────────────────────────────────
\definecolor{darkgray}{gray}{0.15}
\definecolor{medgray}{gray}{0.45}
\definecolor{lightgray}{gray}{0.93}
\definecolor{boxgray}{gray}{0.88}

% ── Section formatting ──────────────────────────────────────────────────────────
\titleformat{\section}[block]
  {\normalfont\large\bfseries}{}
  {0pt}{}[\vspace{2pt}\hrule\vspace{6pt}]
\titlespacing*{\section}{0pt}{14pt}{4pt}

% ── List formatting ─────────────────────────────────────────────────────────────
\setlist[enumerate,1]{label=\textbf{\arabic*.}, leftmargin=2em, itemsep=10pt}

% ── Branding constants (edit here, not per-file) ────────────────────────────────
\newcommand{\SpeedExamLine}{\textbf{Speed Maths:} Competition \& Exam Prep}
\newcommand{\SpeedCredit}{Series by \href{https://www.linkedin.com/in/cerealdev/}{David Akinyele-Aje}}

% ── Header / footer ──────────────────────────────────────────────────────────────
% Usage (once, after % main (standalone preamble for editing)
% vim: nowrap: spell:
% ══════════════════════════════════════════════════════════════════════════════
% Speed Maths --- shared preamble
% Included by every sheet and answer file via:  \input{../../shared/preamble}
% Editing THIS file changes the look of every sheet/answer across every pillar.
% ══════════════════════════════════════════════════════════════════════════════

\usepackage[a4paper, top=25mm, bottom=25mm, left=22mm, right=22mm]{geometry}
\usepackage{amsmath, amssymb}
\usepackage{enumitem}
\usepackage{booktabs}
\usepackage{array}
\usepackage{parskip}
\usepackage{microtype}
\usepackage{fancyhdr}
\usepackage{titlesec}
\usepackage{mdframed}
\usepackage{xcolor}
\usepackage[hidelinks]{hyperref}
\usepackage{graphicx}
\usepackage{tikz}
\usepackage{pgfplots}
\pgfplotsset{compat=1.18}

% ── Colours ───────────────────────────────────────────────────────────────────
\definecolor{darkgray}{gray}{0.15}
\definecolor{medgray}{gray}{0.45}
\definecolor{lightgray}{gray}{0.93}
\definecolor{boxgray}{gray}{0.88}

% ── Section formatting ──────────────────────────────────────────────────────────
\titleformat{\section}[block]
  {\normalfont\large\bfseries}{}
  {0pt}{}[\vspace{2pt}\hrule\vspace{6pt}]
\titlespacing*{\section}{0pt}{14pt}{4pt}

% ── List formatting ─────────────────────────────────────────────────────────────
\setlist[enumerate,1]{label=\textbf{\arabic*.}, leftmargin=2em, itemsep=10pt}

% ── Branding constants (edit here, not per-file) ────────────────────────────────
\newcommand{\SpeedExamLine}{\textbf{Speed Maths:} Competition \& Exam Prep}
\newcommand{\SpeedCredit}{Series by \href{https://www.linkedin.com/in/cerealdev/}{David Akinyele-Aje}}

% ── Header / footer ──────────────────────────────────────────────────────────────
% Usage (once, after \input{../../shared/preamble} and before \begin{document}):
%   \SpeedHeader{Algebra}{3}
% First arg  = pillar name shown as "Speed <Pillar> --- Daily Drill"
% Second arg = worksheet number
\pagestyle{fancy}
\newcommand{\SpeedHeader}[2]{%
  \fancyhf{}%
  \renewcommand{\headrulewidth}{0.4pt}%
  \setlength{\headheight}{14pt}%
  \fancyhead[L]{\small\textbf{Speed #1 --- Daily Drill}}%
  \fancyhead[R]{\small Worksheet \##2}%
  \fancyfoot[L]{\small \SpeedExamLine}%
  \fancyfoot[C]{\small\thepage}%
  \fancyfoot[R]{\small \SpeedCredit}%
  \renewcommand{\footrulewidth}{0.4pt}%
}

% ── Title block (used at the top of every sheet AND answer file) ───────────────
% Usage: \SpeedTitleBlock{Daily Algebra Drill \#3}{Jane Doe}
% %% AGENTS: When generating a new sheet, ask the human contributor for the name they want credited in argument 2.
% For answer files, pass the "--- Answers \& Investigations" suffix in the arg.
\newcommand{\SpeedTitleBlock}[2]{%
  \begin{center}
    {\LARGE\bfseries #1}\\[4pt]
    {\normalsize\itshape Sheet by #2}\\[6pt]
    \hrule\vspace{4pt}
    \begin{tabular}{llcc}
      \toprule
      \textbf{Section} & \textbf{Topic} & \textbf{Questions} & \textbf{Target time}\\
      \midrule
      A & Rapid Recognition          & 10 & 2 min 30 sec \\
      B & Manipulation Drills        & 10 & 8 min        \\
      C & Substitution \& Structure  & 8  & 10 min       \\
      D & Challenge                  & 5  & 15 min       \\
      \bottomrule
    \end{tabular}
    \vspace{4pt}
    \hrule
  \end{center}
}

% ── Sheet metadata: topic + the day's new tools ────────────────────────────────
% Usage (immediately after \SpeedTitleBlock, sheets only):
%   \SpeedMeta{Expanding, factorising, and the identity toolkit}
%             {difference of squares, sum/product form, completing the square}
%
% Arg 1 = the sheet's topic in one line. The website lists this instead of
%         "Sheet 03", so write the thing a reader would actually search for.
% Arg 2 = comma-separated, at most five: the tools this sheet introduces that
%         earlier sheets in the pillar did not. These become the website's tags.
%         Leave out anything the reader already met — the tags are meant to say
%         what is new today, not to summarise the sheet.
%
% Both are parsed straight out of this file by tools/build_website.py, so the
% sheet stays the single source of truth and the site cannot drift from it.
%
% This renders NOTHING. It is a declaration for the build script, not a piece of
% the sheet's layout: adding it to a sheet must not change that sheet's PDF, so
% the 25 existing sheets can be annotated without a single one of them being
% reissued. If the toolkit should also appear on the page, that is a separate
% decision about the sheet design — make it deliberately, here, not by leaving
% this macro printing something nobody asked it to.
\newcommand{\SpeedMeta}[2]{}

% ── Answer / Method / Investigate-further boxes (answer files only) ────────────
\newmdenv[
  backgroundcolor=lightgray,
  linecolor=medgray,
  linewidth=0.5pt,
  leftmargin=0pt, rightmargin=0pt,
  innerleftmargin=8pt, innerrightmargin=8pt,
  innertopmargin=5pt, innerbottommargin=5pt,
  skipabove=4pt, skipbelow=4pt
]{ansbox}

\newmdenv[
  backgroundcolor=white,
  linecolor=darkgray,
  linewidth=0.8pt,
  leftline=true, rightline=false, topline=false, bottomline=false,
  leftmargin=0pt, rightmargin=0pt,
  innerleftmargin=8pt, innerrightmargin=6pt,
  innertopmargin=4pt, innerbottommargin=4pt,
  skipabove=4pt, skipbelow=6pt
]{invbox}

\newcommand{\ans}[1]{%
  \begin{ansbox}
    \textbf{Answer:} #1
  \end{ansbox}}

\newcommand{\method}[1]{%
  {\small\textbf{Method:} #1}\par}

\newcommand{\inv}[1]{%
  \begin{invbox}
    \textbf{\small Investigate further:} {\small #1}
  \end{invbox}}

% ── Misc helpers ────────────────────────────────────────────────────────────────
\newcommand{\qn}[1]{\textbf{#1.}\;}

% Mistake-linking: attach a Top 5 Patterns / Common Traps bullet to the question
% label(s) in this sheet where it actually shows up, e.g. \seealso{B6, D2}
\newcommand{\seealso}[1]{\hfill\textit{\footnotesize(see #1)}}

% ── Graph options macro ─────────────────────────────────────────────────────────
% Usage: \GraphOptions{tikz code for A}{tikz code for B}{tikz code for C}{tikz code for D}
\newcommand{\GraphOptions}[4]{%
  \begin{center}
    \begin{tabular}{cc}
      \textbf{A)} & \textbf{B)} \\
      #1 & #2 \\[10pt]
      \textbf{C)} & \textbf{D)} \\
      #3 & #4
    \end{tabular}
  \end{center}
}

% ── Closing motivational rule + cross-promo footer (sheets only) ────────────────
% Usage: \SpeedClosing{"Quote text."}
\newcommand{\SpeedClosing}[1]{%
  \vspace{20pt}
  \begin{center}
  \rule{0.6\textwidth}{0.4pt}\\[4pt]
  {\small\itshape #1}\\[4pt]
  \rule{0.6\textwidth}{0.4pt}
  \end{center}
  \begin{center}
  {\small Found a mistake or want to contribute a sheet?}\\
  {\small Visit the open-source project at \href{https://github.com/The-CerealDev/speed-maths}{github.com/The-CerealDev/speed-maths}}
  \end{center}
}
 and before \begin{document}):
%   \SpeedHeader{Algebra}{3}
% First arg  = pillar name shown as "Speed <Pillar> --- Daily Drill"
% Second arg = worksheet number
\pagestyle{fancy}
\newcommand{\SpeedHeader}[2]{%
  \fancyhf{}%
  \renewcommand{\headrulewidth}{0.4pt}%
  \setlength{\headheight}{14pt}%
  \fancyhead[L]{\small\textbf{Speed #1 --- Daily Drill}}%
  \fancyhead[R]{\small Worksheet \##2}%
  \fancyfoot[L]{\small \SpeedExamLine}%
  \fancyfoot[C]{\small\thepage}%
  \fancyfoot[R]{\small \SpeedCredit}%
  \renewcommand{\footrulewidth}{0.4pt}%
}

% ── Title block (used at the top of every sheet AND answer file) ───────────────
% Usage: \SpeedTitleBlock{Daily Algebra Drill \#3}{Jane Doe}
% %% AGENTS: When generating a new sheet, ask the human contributor for the name they want credited in argument 2.
% For answer files, pass the "--- Answers \& Investigations" suffix in the arg.
\newcommand{\SpeedTitleBlock}[2]{%
  \begin{center}
    {\LARGE\bfseries #1}\\[4pt]
    {\normalsize\itshape Sheet by #2}\\[6pt]
    \hrule\vspace{4pt}
    \begin{tabular}{llcc}
      \toprule
      \textbf{Section} & \textbf{Topic} & \textbf{Questions} & \textbf{Target time}\\
      \midrule
      A & Rapid Recognition          & 10 & 2 min 30 sec \\
      B & Manipulation Drills        & 10 & 8 min        \\
      C & Substitution \& Structure  & 8  & 10 min       \\
      D & Challenge                  & 5  & 15 min       \\
      \bottomrule
    \end{tabular}
    \vspace{4pt}
    \hrule
  \end{center}
}

% ── Sheet metadata: topic + the day's new tools ────────────────────────────────
% Usage (immediately after \SpeedTitleBlock, sheets only):
%   \SpeedMeta{Expanding, factorising, and the identity toolkit}
%             {difference of squares, sum/product form, completing the square}
%
% Arg 1 = the sheet's topic in one line. The website lists this instead of
%         "Sheet 03", so write the thing a reader would actually search for.
% Arg 2 = comma-separated, at most five: the tools this sheet introduces that
%         earlier sheets in the pillar did not. These become the website's tags.
%         Leave out anything the reader already met — the tags are meant to say
%         what is new today, not to summarise the sheet.
%
% Both are parsed straight out of this file by tools/build_website.py, so the
% sheet stays the single source of truth and the site cannot drift from it.
%
% This renders NOTHING. It is a declaration for the build script, not a piece of
% the sheet's layout: adding it to a sheet must not change that sheet's PDF, so
% the 25 existing sheets can be annotated without a single one of them being
% reissued. If the toolkit should also appear on the page, that is a separate
% decision about the sheet design — make it deliberately, here, not by leaving
% this macro printing something nobody asked it to.
\newcommand{\SpeedMeta}[2]{}

% ── Answer / Method / Investigate-further boxes (answer files only) ────────────
\newmdenv[
  backgroundcolor=lightgray,
  linecolor=medgray,
  linewidth=0.5pt,
  leftmargin=0pt, rightmargin=0pt,
  innerleftmargin=8pt, innerrightmargin=8pt,
  innertopmargin=5pt, innerbottommargin=5pt,
  skipabove=4pt, skipbelow=4pt
]{ansbox}

\newmdenv[
  backgroundcolor=white,
  linecolor=darkgray,
  linewidth=0.8pt,
  leftline=true, rightline=false, topline=false, bottomline=false,
  leftmargin=0pt, rightmargin=0pt,
  innerleftmargin=8pt, innerrightmargin=6pt,
  innertopmargin=4pt, innerbottommargin=4pt,
  skipabove=4pt, skipbelow=6pt
]{invbox}

\newcommand{\ans}[1]{%
  \begin{ansbox}
    \textbf{Answer:} #1
  \end{ansbox}}

\newcommand{\method}[1]{%
  {\small\textbf{Method:} #1}\par}

\newcommand{\inv}[1]{%
  \begin{invbox}
    \textbf{\small Investigate further:} {\small #1}
  \end{invbox}}

% ── Misc helpers ────────────────────────────────────────────────────────────────
\newcommand{\qn}[1]{\textbf{#1.}\;}

% Mistake-linking: attach a Top 5 Patterns / Common Traps bullet to the question
% label(s) in this sheet where it actually shows up, e.g. \seealso{B6, D2}
\newcommand{\seealso}[1]{\hfill\textit{\footnotesize(see #1)}}

% ── Graph options macro ─────────────────────────────────────────────────────────
% Usage: \GraphOptions{tikz code for A}{tikz code for B}{tikz code for C}{tikz code for D}
\newcommand{\GraphOptions}[4]{%
  \begin{center}
    \begin{tabular}{cc}
      \textbf{A)} & \textbf{B)} \\
      #1 & #2 \\[10pt]
      \textbf{C)} & \textbf{D)} \\
      #3 & #4
    \end{tabular}
  \end{center}
}

% ── Closing motivational rule + cross-promo footer (sheets only) ────────────────
% Usage: \SpeedClosing{"Quote text."}
\newcommand{\SpeedClosing}[1]{%
  \vspace{20pt}
  \begin{center}
  \rule{0.6\textwidth}{0.4pt}\\[4pt]
  {\small\itshape #1}\\[4pt]
  \rule{0.6\textwidth}{0.4pt}
  \end{center}
  \begin{center}
  {\small Found a mistake or want to contribute a sheet?}\\
  {\small Visit the open-source project at \href{https://github.com/The-CerealDev/speed-maths}{github.com/The-CerealDev/speed-maths}}
  \end{center}
}

% Editing THIS file changes the look of every sheet/answer across every pillar.
% ══════════════════════════════════════════════════════════════════════════════

\usepackage[a4paper, top=25mm, bottom=25mm, left=22mm, right=22mm]{geometry}
\usepackage{amsmath, amssymb}
\usepackage{enumitem}
\usepackage{booktabs}
\usepackage{array}
\usepackage{parskip}
\usepackage{microtype}
\usepackage{fancyhdr}
\usepackage{titlesec}
\usepackage{mdframed}
\usepackage{xcolor}
\usepackage[hidelinks]{hyperref}
\usepackage{graphicx}
\usepackage{tikz}
\usepackage{pgfplots}
\pgfplotsset{compat=1.18}

% ── Colours ───────────────────────────────────────────────────────────────────
\definecolor{darkgray}{gray}{0.15}
\definecolor{medgray}{gray}{0.45}
\definecolor{lightgray}{gray}{0.93}
\definecolor{boxgray}{gray}{0.88}

% ── Section formatting ──────────────────────────────────────────────────────────
\titleformat{\section}[block]
  {\normalfont\large\bfseries}{}
  {0pt}{}[\vspace{2pt}\hrule\vspace{6pt}]
\titlespacing*{\section}{0pt}{14pt}{4pt}

% ── List formatting ─────────────────────────────────────────────────────────────
\setlist[enumerate,1]{label=\textbf{\arabic*.}, leftmargin=2em, itemsep=10pt}

% ── Branding constants (edit here, not per-file) ────────────────────────────────
\newcommand{\SpeedExamLine}{\textbf{Speed Maths:} Competition \& Exam Prep}
\newcommand{\SpeedCredit}{Series by \href{https://www.linkedin.com/in/cerealdev/}{David Akinyele-Aje}}

% ── Header / footer ──────────────────────────────────────────────────────────────
% Usage (once, after % main (standalone preamble for editing)
% vim: nowrap: spell:
% ══════════════════════════════════════════════════════════════════════════════
% Speed Maths --- shared preamble
% Included by every sheet and answer file via:  % main (standalone preamble for editing)
% vim: nowrap: spell:
% ══════════════════════════════════════════════════════════════════════════════
% Speed Maths --- shared preamble
% Included by every sheet and answer file via:  \input{../../shared/preamble}
% Editing THIS file changes the look of every sheet/answer across every pillar.
% ══════════════════════════════════════════════════════════════════════════════

\usepackage[a4paper, top=25mm, bottom=25mm, left=22mm, right=22mm]{geometry}
\usepackage{amsmath, amssymb}
\usepackage{enumitem}
\usepackage{booktabs}
\usepackage{array}
\usepackage{parskip}
\usepackage{microtype}
\usepackage{fancyhdr}
\usepackage{titlesec}
\usepackage{mdframed}
\usepackage{xcolor}
\usepackage[hidelinks]{hyperref}
\usepackage{graphicx}
\usepackage{tikz}
\usepackage{pgfplots}
\pgfplotsset{compat=1.18}

% ── Colours ───────────────────────────────────────────────────────────────────
\definecolor{darkgray}{gray}{0.15}
\definecolor{medgray}{gray}{0.45}
\definecolor{lightgray}{gray}{0.93}
\definecolor{boxgray}{gray}{0.88}

% ── Section formatting ──────────────────────────────────────────────────────────
\titleformat{\section}[block]
  {\normalfont\large\bfseries}{}
  {0pt}{}[\vspace{2pt}\hrule\vspace{6pt}]
\titlespacing*{\section}{0pt}{14pt}{4pt}

% ── List formatting ─────────────────────────────────────────────────────────────
\setlist[enumerate,1]{label=\textbf{\arabic*.}, leftmargin=2em, itemsep=10pt}

% ── Branding constants (edit here, not per-file) ────────────────────────────────
\newcommand{\SpeedExamLine}{\textbf{Speed Maths:} Competition \& Exam Prep}
\newcommand{\SpeedCredit}{Series by \href{https://www.linkedin.com/in/cerealdev/}{David Akinyele-Aje}}

% ── Header / footer ──────────────────────────────────────────────────────────────
% Usage (once, after \input{../../shared/preamble} and before \begin{document}):
%   \SpeedHeader{Algebra}{3}
% First arg  = pillar name shown as "Speed <Pillar> --- Daily Drill"
% Second arg = worksheet number
\pagestyle{fancy}
\newcommand{\SpeedHeader}[2]{%
  \fancyhf{}%
  \renewcommand{\headrulewidth}{0.4pt}%
  \setlength{\headheight}{14pt}%
  \fancyhead[L]{\small\textbf{Speed #1 --- Daily Drill}}%
  \fancyhead[R]{\small Worksheet \##2}%
  \fancyfoot[L]{\small \SpeedExamLine}%
  \fancyfoot[C]{\small\thepage}%
  \fancyfoot[R]{\small \SpeedCredit}%
  \renewcommand{\footrulewidth}{0.4pt}%
}

% ── Title block (used at the top of every sheet AND answer file) ───────────────
% Usage: \SpeedTitleBlock{Daily Algebra Drill \#3}{Jane Doe}
% %% AGENTS: When generating a new sheet, ask the human contributor for the name they want credited in argument 2.
% For answer files, pass the "--- Answers \& Investigations" suffix in the arg.
\newcommand{\SpeedTitleBlock}[2]{%
  \begin{center}
    {\LARGE\bfseries #1}\\[4pt]
    {\normalsize\itshape Sheet by #2}\\[6pt]
    \hrule\vspace{4pt}
    \begin{tabular}{llcc}
      \toprule
      \textbf{Section} & \textbf{Topic} & \textbf{Questions} & \textbf{Target time}\\
      \midrule
      A & Rapid Recognition          & 10 & 2 min 30 sec \\
      B & Manipulation Drills        & 10 & 8 min        \\
      C & Substitution \& Structure  & 8  & 10 min       \\
      D & Challenge                  & 5  & 15 min       \\
      \bottomrule
    \end{tabular}
    \vspace{4pt}
    \hrule
  \end{center}
}

% ── Sheet metadata: topic + the day's new tools ────────────────────────────────
% Usage (immediately after \SpeedTitleBlock, sheets only):
%   \SpeedMeta{Expanding, factorising, and the identity toolkit}
%             {difference of squares, sum/product form, completing the square}
%
% Arg 1 = the sheet's topic in one line. The website lists this instead of
%         "Sheet 03", so write the thing a reader would actually search for.
% Arg 2 = comma-separated, at most five: the tools this sheet introduces that
%         earlier sheets in the pillar did not. These become the website's tags.
%         Leave out anything the reader already met — the tags are meant to say
%         what is new today, not to summarise the sheet.
%
% Both are parsed straight out of this file by tools/build_website.py, so the
% sheet stays the single source of truth and the site cannot drift from it.
%
% This renders NOTHING. It is a declaration for the build script, not a piece of
% the sheet's layout: adding it to a sheet must not change that sheet's PDF, so
% the 25 existing sheets can be annotated without a single one of them being
% reissued. If the toolkit should also appear on the page, that is a separate
% decision about the sheet design — make it deliberately, here, not by leaving
% this macro printing something nobody asked it to.
\newcommand{\SpeedMeta}[2]{}

% ── Answer / Method / Investigate-further boxes (answer files only) ────────────
\newmdenv[
  backgroundcolor=lightgray,
  linecolor=medgray,
  linewidth=0.5pt,
  leftmargin=0pt, rightmargin=0pt,
  innerleftmargin=8pt, innerrightmargin=8pt,
  innertopmargin=5pt, innerbottommargin=5pt,
  skipabove=4pt, skipbelow=4pt
]{ansbox}

\newmdenv[
  backgroundcolor=white,
  linecolor=darkgray,
  linewidth=0.8pt,
  leftline=true, rightline=false, topline=false, bottomline=false,
  leftmargin=0pt, rightmargin=0pt,
  innerleftmargin=8pt, innerrightmargin=6pt,
  innertopmargin=4pt, innerbottommargin=4pt,
  skipabove=4pt, skipbelow=6pt
]{invbox}

\newcommand{\ans}[1]{%
  \begin{ansbox}
    \textbf{Answer:} #1
  \end{ansbox}}

\newcommand{\method}[1]{%
  {\small\textbf{Method:} #1}\par}

\newcommand{\inv}[1]{%
  \begin{invbox}
    \textbf{\small Investigate further:} {\small #1}
  \end{invbox}}

% ── Misc helpers ────────────────────────────────────────────────────────────────
\newcommand{\qn}[1]{\textbf{#1.}\;}

% Mistake-linking: attach a Top 5 Patterns / Common Traps bullet to the question
% label(s) in this sheet where it actually shows up, e.g. \seealso{B6, D2}
\newcommand{\seealso}[1]{\hfill\textit{\footnotesize(see #1)}}

% ── Graph options macro ─────────────────────────────────────────────────────────
% Usage: \GraphOptions{tikz code for A}{tikz code for B}{tikz code for C}{tikz code for D}
\newcommand{\GraphOptions}[4]{%
  \begin{center}
    \begin{tabular}{cc}
      \textbf{A)} & \textbf{B)} \\
      #1 & #2 \\[10pt]
      \textbf{C)} & \textbf{D)} \\
      #3 & #4
    \end{tabular}
  \end{center}
}

% ── Closing motivational rule + cross-promo footer (sheets only) ────────────────
% Usage: \SpeedClosing{"Quote text."}
\newcommand{\SpeedClosing}[1]{%
  \vspace{20pt}
  \begin{center}
  \rule{0.6\textwidth}{0.4pt}\\[4pt]
  {\small\itshape #1}\\[4pt]
  \rule{0.6\textwidth}{0.4pt}
  \end{center}
  \begin{center}
  {\small Found a mistake or want to contribute a sheet?}\\
  {\small Visit the open-source project at \href{https://github.com/The-CerealDev/speed-maths}{github.com/The-CerealDev/speed-maths}}
  \end{center}
}

% Editing THIS file changes the look of every sheet/answer across every pillar.
% ══════════════════════════════════════════════════════════════════════════════

\usepackage[a4paper, top=25mm, bottom=25mm, left=22mm, right=22mm]{geometry}
\usepackage{amsmath, amssymb}
\usepackage{enumitem}
\usepackage{booktabs}
\usepackage{array}
\usepackage{parskip}
\usepackage{microtype}
\usepackage{fancyhdr}
\usepackage{titlesec}
\usepackage{mdframed}
\usepackage{xcolor}
\usepackage[hidelinks]{hyperref}
\usepackage{graphicx}
\usepackage{tikz}
\usepackage{pgfplots}
\pgfplotsset{compat=1.18}

% ── Colours ───────────────────────────────────────────────────────────────────
\definecolor{darkgray}{gray}{0.15}
\definecolor{medgray}{gray}{0.45}
\definecolor{lightgray}{gray}{0.93}
\definecolor{boxgray}{gray}{0.88}

% ── Section formatting ──────────────────────────────────────────────────────────
\titleformat{\section}[block]
  {\normalfont\large\bfseries}{}
  {0pt}{}[\vspace{2pt}\hrule\vspace{6pt}]
\titlespacing*{\section}{0pt}{14pt}{4pt}

% ── List formatting ─────────────────────────────────────────────────────────────
\setlist[enumerate,1]{label=\textbf{\arabic*.}, leftmargin=2em, itemsep=10pt}

% ── Branding constants (edit here, not per-file) ────────────────────────────────
\newcommand{\SpeedExamLine}{\textbf{Speed Maths:} Competition \& Exam Prep}
\newcommand{\SpeedCredit}{Series by \href{https://www.linkedin.com/in/cerealdev/}{David Akinyele-Aje}}

% ── Header / footer ──────────────────────────────────────────────────────────────
% Usage (once, after % main (standalone preamble for editing)
% vim: nowrap: spell:
% ══════════════════════════════════════════════════════════════════════════════
% Speed Maths --- shared preamble
% Included by every sheet and answer file via:  \input{../../shared/preamble}
% Editing THIS file changes the look of every sheet/answer across every pillar.
% ══════════════════════════════════════════════════════════════════════════════

\usepackage[a4paper, top=25mm, bottom=25mm, left=22mm, right=22mm]{geometry}
\usepackage{amsmath, amssymb}
\usepackage{enumitem}
\usepackage{booktabs}
\usepackage{array}
\usepackage{parskip}
\usepackage{microtype}
\usepackage{fancyhdr}
\usepackage{titlesec}
\usepackage{mdframed}
\usepackage{xcolor}
\usepackage[hidelinks]{hyperref}
\usepackage{graphicx}
\usepackage{tikz}
\usepackage{pgfplots}
\pgfplotsset{compat=1.18}

% ── Colours ───────────────────────────────────────────────────────────────────
\definecolor{darkgray}{gray}{0.15}
\definecolor{medgray}{gray}{0.45}
\definecolor{lightgray}{gray}{0.93}
\definecolor{boxgray}{gray}{0.88}

% ── Section formatting ──────────────────────────────────────────────────────────
\titleformat{\section}[block]
  {\normalfont\large\bfseries}{}
  {0pt}{}[\vspace{2pt}\hrule\vspace{6pt}]
\titlespacing*{\section}{0pt}{14pt}{4pt}

% ── List formatting ─────────────────────────────────────────────────────────────
\setlist[enumerate,1]{label=\textbf{\arabic*.}, leftmargin=2em, itemsep=10pt}

% ── Branding constants (edit here, not per-file) ────────────────────────────────
\newcommand{\SpeedExamLine}{\textbf{Speed Maths:} Competition \& Exam Prep}
\newcommand{\SpeedCredit}{Series by \href{https://www.linkedin.com/in/cerealdev/}{David Akinyele-Aje}}

% ── Header / footer ──────────────────────────────────────────────────────────────
% Usage (once, after \input{../../shared/preamble} and before \begin{document}):
%   \SpeedHeader{Algebra}{3}
% First arg  = pillar name shown as "Speed <Pillar> --- Daily Drill"
% Second arg = worksheet number
\pagestyle{fancy}
\newcommand{\SpeedHeader}[2]{%
  \fancyhf{}%
  \renewcommand{\headrulewidth}{0.4pt}%
  \setlength{\headheight}{14pt}%
  \fancyhead[L]{\small\textbf{Speed #1 --- Daily Drill}}%
  \fancyhead[R]{\small Worksheet \##2}%
  \fancyfoot[L]{\small \SpeedExamLine}%
  \fancyfoot[C]{\small\thepage}%
  \fancyfoot[R]{\small \SpeedCredit}%
  \renewcommand{\footrulewidth}{0.4pt}%
}

% ── Title block (used at the top of every sheet AND answer file) ───────────────
% Usage: \SpeedTitleBlock{Daily Algebra Drill \#3}{Jane Doe}
% %% AGENTS: When generating a new sheet, ask the human contributor for the name they want credited in argument 2.
% For answer files, pass the "--- Answers \& Investigations" suffix in the arg.
\newcommand{\SpeedTitleBlock}[2]{%
  \begin{center}
    {\LARGE\bfseries #1}\\[4pt]
    {\normalsize\itshape Sheet by #2}\\[6pt]
    \hrule\vspace{4pt}
    \begin{tabular}{llcc}
      \toprule
      \textbf{Section} & \textbf{Topic} & \textbf{Questions} & \textbf{Target time}\\
      \midrule
      A & Rapid Recognition          & 10 & 2 min 30 sec \\
      B & Manipulation Drills        & 10 & 8 min        \\
      C & Substitution \& Structure  & 8  & 10 min       \\
      D & Challenge                  & 5  & 15 min       \\
      \bottomrule
    \end{tabular}
    \vspace{4pt}
    \hrule
  \end{center}
}

% ── Sheet metadata: topic + the day's new tools ────────────────────────────────
% Usage (immediately after \SpeedTitleBlock, sheets only):
%   \SpeedMeta{Expanding, factorising, and the identity toolkit}
%             {difference of squares, sum/product form, completing the square}
%
% Arg 1 = the sheet's topic in one line. The website lists this instead of
%         "Sheet 03", so write the thing a reader would actually search for.
% Arg 2 = comma-separated, at most five: the tools this sheet introduces that
%         earlier sheets in the pillar did not. These become the website's tags.
%         Leave out anything the reader already met — the tags are meant to say
%         what is new today, not to summarise the sheet.
%
% Both are parsed straight out of this file by tools/build_website.py, so the
% sheet stays the single source of truth and the site cannot drift from it.
%
% This renders NOTHING. It is a declaration for the build script, not a piece of
% the sheet's layout: adding it to a sheet must not change that sheet's PDF, so
% the 25 existing sheets can be annotated without a single one of them being
% reissued. If the toolkit should also appear on the page, that is a separate
% decision about the sheet design — make it deliberately, here, not by leaving
% this macro printing something nobody asked it to.
\newcommand{\SpeedMeta}[2]{}

% ── Answer / Method / Investigate-further boxes (answer files only) ────────────
\newmdenv[
  backgroundcolor=lightgray,
  linecolor=medgray,
  linewidth=0.5pt,
  leftmargin=0pt, rightmargin=0pt,
  innerleftmargin=8pt, innerrightmargin=8pt,
  innertopmargin=5pt, innerbottommargin=5pt,
  skipabove=4pt, skipbelow=4pt
]{ansbox}

\newmdenv[
  backgroundcolor=white,
  linecolor=darkgray,
  linewidth=0.8pt,
  leftline=true, rightline=false, topline=false, bottomline=false,
  leftmargin=0pt, rightmargin=0pt,
  innerleftmargin=8pt, innerrightmargin=6pt,
  innertopmargin=4pt, innerbottommargin=4pt,
  skipabove=4pt, skipbelow=6pt
]{invbox}

\newcommand{\ans}[1]{%
  \begin{ansbox}
    \textbf{Answer:} #1
  \end{ansbox}}

\newcommand{\method}[1]{%
  {\small\textbf{Method:} #1}\par}

\newcommand{\inv}[1]{%
  \begin{invbox}
    \textbf{\small Investigate further:} {\small #1}
  \end{invbox}}

% ── Misc helpers ────────────────────────────────────────────────────────────────
\newcommand{\qn}[1]{\textbf{#1.}\;}

% Mistake-linking: attach a Top 5 Patterns / Common Traps bullet to the question
% label(s) in this sheet where it actually shows up, e.g. \seealso{B6, D2}
\newcommand{\seealso}[1]{\hfill\textit{\footnotesize(see #1)}}

% ── Graph options macro ─────────────────────────────────────────────────────────
% Usage: \GraphOptions{tikz code for A}{tikz code for B}{tikz code for C}{tikz code for D}
\newcommand{\GraphOptions}[4]{%
  \begin{center}
    \begin{tabular}{cc}
      \textbf{A)} & \textbf{B)} \\
      #1 & #2 \\[10pt]
      \textbf{C)} & \textbf{D)} \\
      #3 & #4
    \end{tabular}
  \end{center}
}

% ── Closing motivational rule + cross-promo footer (sheets only) ────────────────
% Usage: \SpeedClosing{"Quote text."}
\newcommand{\SpeedClosing}[1]{%
  \vspace{20pt}
  \begin{center}
  \rule{0.6\textwidth}{0.4pt}\\[4pt]
  {\small\itshape #1}\\[4pt]
  \rule{0.6\textwidth}{0.4pt}
  \end{center}
  \begin{center}
  {\small Found a mistake or want to contribute a sheet?}\\
  {\small Visit the open-source project at \href{https://github.com/The-CerealDev/speed-maths}{github.com/The-CerealDev/speed-maths}}
  \end{center}
}
 and before \begin{document}):
%   \SpeedHeader{Algebra}{3}
% First arg  = pillar name shown as "Speed <Pillar> --- Daily Drill"
% Second arg = worksheet number
\pagestyle{fancy}
\newcommand{\SpeedHeader}[2]{%
  \fancyhf{}%
  \renewcommand{\headrulewidth}{0.4pt}%
  \setlength{\headheight}{14pt}%
  \fancyhead[L]{\small\textbf{Speed #1 --- Daily Drill}}%
  \fancyhead[R]{\small Worksheet \##2}%
  \fancyfoot[L]{\small \SpeedExamLine}%
  \fancyfoot[C]{\small\thepage}%
  \fancyfoot[R]{\small \SpeedCredit}%
  \renewcommand{\footrulewidth}{0.4pt}%
}

% ── Title block (used at the top of every sheet AND answer file) ───────────────
% Usage: \SpeedTitleBlock{Daily Algebra Drill \#3}{Jane Doe}
% %% AGENTS: When generating a new sheet, ask the human contributor for the name they want credited in argument 2.
% For answer files, pass the "--- Answers \& Investigations" suffix in the arg.
\newcommand{\SpeedTitleBlock}[2]{%
  \begin{center}
    {\LARGE\bfseries #1}\\[4pt]
    {\normalsize\itshape Sheet by #2}\\[6pt]
    \hrule\vspace{4pt}
    \begin{tabular}{llcc}
      \toprule
      \textbf{Section} & \textbf{Topic} & \textbf{Questions} & \textbf{Target time}\\
      \midrule
      A & Rapid Recognition          & 10 & 2 min 30 sec \\
      B & Manipulation Drills        & 10 & 8 min        \\
      C & Substitution \& Structure  & 8  & 10 min       \\
      D & Challenge                  & 5  & 15 min       \\
      \bottomrule
    \end{tabular}
    \vspace{4pt}
    \hrule
  \end{center}
}

% ── Sheet metadata: topic + the day's new tools ────────────────────────────────
% Usage (immediately after \SpeedTitleBlock, sheets only):
%   \SpeedMeta{Expanding, factorising, and the identity toolkit}
%             {difference of squares, sum/product form, completing the square}
%
% Arg 1 = the sheet's topic in one line. The website lists this instead of
%         "Sheet 03", so write the thing a reader would actually search for.
% Arg 2 = comma-separated, at most five: the tools this sheet introduces that
%         earlier sheets in the pillar did not. These become the website's tags.
%         Leave out anything the reader already met — the tags are meant to say
%         what is new today, not to summarise the sheet.
%
% Both are parsed straight out of this file by tools/build_website.py, so the
% sheet stays the single source of truth and the site cannot drift from it.
%
% This renders NOTHING. It is a declaration for the build script, not a piece of
% the sheet's layout: adding it to a sheet must not change that sheet's PDF, so
% the 25 existing sheets can be annotated without a single one of them being
% reissued. If the toolkit should also appear on the page, that is a separate
% decision about the sheet design — make it deliberately, here, not by leaving
% this macro printing something nobody asked it to.
\newcommand{\SpeedMeta}[2]{}

% ── Answer / Method / Investigate-further boxes (answer files only) ────────────
\newmdenv[
  backgroundcolor=lightgray,
  linecolor=medgray,
  linewidth=0.5pt,
  leftmargin=0pt, rightmargin=0pt,
  innerleftmargin=8pt, innerrightmargin=8pt,
  innertopmargin=5pt, innerbottommargin=5pt,
  skipabove=4pt, skipbelow=4pt
]{ansbox}

\newmdenv[
  backgroundcolor=white,
  linecolor=darkgray,
  linewidth=0.8pt,
  leftline=true, rightline=false, topline=false, bottomline=false,
  leftmargin=0pt, rightmargin=0pt,
  innerleftmargin=8pt, innerrightmargin=6pt,
  innertopmargin=4pt, innerbottommargin=4pt,
  skipabove=4pt, skipbelow=6pt
]{invbox}

\newcommand{\ans}[1]{%
  \begin{ansbox}
    \textbf{Answer:} #1
  \end{ansbox}}

\newcommand{\method}[1]{%
  {\small\textbf{Method:} #1}\par}

\newcommand{\inv}[1]{%
  \begin{invbox}
    \textbf{\small Investigate further:} {\small #1}
  \end{invbox}}

% ── Misc helpers ────────────────────────────────────────────────────────────────
\newcommand{\qn}[1]{\textbf{#1.}\;}

% Mistake-linking: attach a Top 5 Patterns / Common Traps bullet to the question
% label(s) in this sheet where it actually shows up, e.g. \seealso{B6, D2}
\newcommand{\seealso}[1]{\hfill\textit{\footnotesize(see #1)}}

% ── Graph options macro ─────────────────────────────────────────────────────────
% Usage: \GraphOptions{tikz code for A}{tikz code for B}{tikz code for C}{tikz code for D}
\newcommand{\GraphOptions}[4]{%
  \begin{center}
    \begin{tabular}{cc}
      \textbf{A)} & \textbf{B)} \\
      #1 & #2 \\[10pt]
      \textbf{C)} & \textbf{D)} \\
      #3 & #4
    \end{tabular}
  \end{center}
}

% ── Closing motivational rule + cross-promo footer (sheets only) ────────────────
% Usage: \SpeedClosing{"Quote text."}
\newcommand{\SpeedClosing}[1]{%
  \vspace{20pt}
  \begin{center}
  \rule{0.6\textwidth}{0.4pt}\\[4pt]
  {\small\itshape #1}\\[4pt]
  \rule{0.6\textwidth}{0.4pt}
  \end{center}
  \begin{center}
  {\small Found a mistake or want to contribute a sheet?}\\
  {\small Visit the open-source project at \href{https://github.com/The-CerealDev/speed-maths}{github.com/The-CerealDev/speed-maths}}
  \end{center}
}
 and before \begin{document}):
%   \SpeedHeader{Algebra}{3}
% First arg  = pillar name shown as "Speed <Pillar> --- Daily Drill"
% Second arg = worksheet number
\pagestyle{fancy}
\newcommand{\SpeedHeader}[2]{%
  \fancyhf{}%
  \renewcommand{\headrulewidth}{0.4pt}%
  \setlength{\headheight}{14pt}%
  \fancyhead[L]{\small\textbf{Speed #1 --- Daily Drill}}%
  \fancyhead[R]{\small Worksheet \##2}%
  \fancyfoot[L]{\small \SpeedExamLine}%
  \fancyfoot[C]{\small\thepage}%
  \fancyfoot[R]{\small \SpeedCredit}%
  \renewcommand{\footrulewidth}{0.4pt}%
}

% ── Title block (used at the top of every sheet AND answer file) ───────────────
% Usage: \SpeedTitleBlock{Daily Algebra Drill \#3}{Jane Doe}
% %% AGENTS: When generating a new sheet, ask the human contributor for the name they want credited in argument 2.
% For answer files, pass the "--- Answers \& Investigations" suffix in the arg.
\newcommand{\SpeedTitleBlock}[2]{%
  \begin{center}
    {\LARGE\bfseries #1}\\[4pt]
    {\normalsize\itshape Sheet by #2}\\[6pt]
    \hrule\vspace{4pt}
    \begin{tabular}{llcc}
      \toprule
      \textbf{Section} & \textbf{Topic} & \textbf{Questions} & \textbf{Target time}\\
      \midrule
      A & Rapid Recognition          & 10 & 2 min 30 sec \\
      B & Manipulation Drills        & 10 & 8 min        \\
      C & Substitution \& Structure  & 8  & 10 min       \\
      D & Challenge                  & 5  & 15 min       \\
      \bottomrule
    \end{tabular}
    \vspace{4pt}
    \hrule
  \end{center}
}

% ── Sheet metadata: topic + the day's new tools ────────────────────────────────
% Usage (immediately after \SpeedTitleBlock, sheets only):
%   \SpeedMeta{Expanding, factorising, and the identity toolkit}
%             {difference of squares, sum/product form, completing the square}
%
% Arg 1 = the sheet's topic in one line. The website lists this instead of
%         "Sheet 03", so write the thing a reader would actually search for.
% Arg 2 = comma-separated, at most five: the tools this sheet introduces that
%         earlier sheets in the pillar did not. These become the website's tags.
%         Leave out anything the reader already met — the tags are meant to say
%         what is new today, not to summarise the sheet.
%
% Both are parsed straight out of this file by tools/build_website.py, so the
% sheet stays the single source of truth and the site cannot drift from it.
%
% This renders NOTHING. It is a declaration for the build script, not a piece of
% the sheet's layout: adding it to a sheet must not change that sheet's PDF, so
% the 25 existing sheets can be annotated without a single one of them being
% reissued. If the toolkit should also appear on the page, that is a separate
% decision about the sheet design — make it deliberately, here, not by leaving
% this macro printing something nobody asked it to.
\newcommand{\SpeedMeta}[2]{}

% ── Answer / Method / Investigate-further boxes (answer files only) ────────────
\newmdenv[
  backgroundcolor=lightgray,
  linecolor=medgray,
  linewidth=0.5pt,
  leftmargin=0pt, rightmargin=0pt,
  innerleftmargin=8pt, innerrightmargin=8pt,
  innertopmargin=5pt, innerbottommargin=5pt,
  skipabove=4pt, skipbelow=4pt
]{ansbox}

\newmdenv[
  backgroundcolor=white,
  linecolor=darkgray,
  linewidth=0.8pt,
  leftline=true, rightline=false, topline=false, bottomline=false,
  leftmargin=0pt, rightmargin=0pt,
  innerleftmargin=8pt, innerrightmargin=6pt,
  innertopmargin=4pt, innerbottommargin=4pt,
  skipabove=4pt, skipbelow=6pt
]{invbox}

\newcommand{\ans}[1]{%
  \begin{ansbox}
    \textbf{Answer:} #1
  \end{ansbox}}

\newcommand{\method}[1]{%
  {\small\textbf{Method:} #1}\par}

\newcommand{\inv}[1]{%
  \begin{invbox}
    \textbf{\small Investigate further:} {\small #1}
  \end{invbox}}

% ── Misc helpers ────────────────────────────────────────────────────────────────
\newcommand{\qn}[1]{\textbf{#1.}\;}

% Mistake-linking: attach a Top 5 Patterns / Common Traps bullet to the question
% label(s) in this sheet where it actually shows up, e.g. \seealso{B6, D2}
\newcommand{\seealso}[1]{\hfill\textit{\footnotesize(see #1)}}

% ── Graph options macro ─────────────────────────────────────────────────────────
% Usage: \GraphOptions{tikz code for A}{tikz code for B}{tikz code for C}{tikz code for D}
\newcommand{\GraphOptions}[4]{%
  \begin{center}
    \begin{tabular}{cc}
      \textbf{A)} & \textbf{B)} \\
      #1 & #2 \\[10pt]
      \textbf{C)} & \textbf{D)} \\
      #3 & #4
    \end{tabular}
  \end{center}
}

% ── Closing motivational rule + cross-promo footer (sheets only) ────────────────
% Usage: \SpeedClosing{"Quote text."}
\newcommand{\SpeedClosing}[1]{%
  \vspace{20pt}
  \begin{center}
  \rule{0.6\textwidth}{0.4pt}\\[4pt]
  {\small\itshape #1}\\[4pt]
  \rule{0.6\textwidth}{0.4pt}
  \end{center}
  \begin{center}
  {\small Found a mistake or want to contribute a sheet?}\\
  {\small Visit the open-source project at \href{https://github.com/The-CerealDev/speed-maths}{github.com/The-CerealDev/speed-maths}}
  \end{center}
}

\SpeedHeader{Geometry}{7}

\begin{document}

\SpeedTitleBlock{Daily Geometry Drill \#7 --- Answers \& Investigations}{Henry Koduthore}
\SpeedMeta{TMUA-style synthesis: loci and geometric logic}{loci of fixed distance, Apollonius ratio circles, necessary/sufficient conditionals, spot-the-flaw, mixed synthesis}

\section*{Section A \quad Rapid Recognition}

\begin{enumerate}

\item \ans{"a circle of radius 5 centred at the origin"}
\method{$\sqrt{x^2+y^2}=5$ rearranges to $x^2+y^2=25$, a circle of radius $5$.}
\inv{Forgetting to square: quoting radius $25$ or a square.}

\item \ans{$x=2$}
\method{Squaring $x^2+y^2=(x-4)^2+y^2$: $x^2=(x-4)^2$, so $x=2$.}
\inv{Giving the single midpoint $(2,0)$ instead of the whole vertical line.}

\item \ans{$4$}
\method{On the $x$-axis, between $A$ and $B$ we have $PA=x$ and $PB=6-x$, so $x=2(6-x)$, giving $x=4$. The Apollonius circle's two $x$-axis endpoints are $x=4$ and $x=12$.}
\inv{Choosing $x=2$ (halving $6$) or the external endpoint $x=12$.}

\item \ans{$y=1$ or $y=7$}
\method{$|y-4|=3$ gives $y=1$ and $y=7$, two parallel lines $3$ units from $y=4$.}
\inv{Keeping only one line $y=1$ and missing its mirror $y=7$.}

\item \ans{$(2,1)$}
\method{A point equidistant from focus $(0,1)$ and directrix $y=-1$ satisfies $\sqrt{x^2+(y-1)^2}=|y+1|$, i.e.\ $x^2=4y$; $(2,1)$ works since $4=4$.}
\inv{Naming a point on the $y$-axis (which only hits the vertex origin for integer $y$) or mislaying the focus/directrix distance.}

\item \ans{$(x-1)^2+(y-2)^2=9$}
\method{Points at distance exactly $3$ from $(1,2)$ form the circle centred there with radius $3$.}
\inv{Writing distance $=3$ unsquared, or using the point as an offset in the wrong sign.}

\item \ans{"a circle of radius 1 centred at the origin"}
\method{$Q$ runs over $x^2+y^2=4$ (radius $2$); the midpoint $M=OQ/2$ scales radii by $\frac12$, giving $x^2+y^2=1$.}
\inv{Halving the equation's radius linearly to $\sqrt{2}$ or keeping the original radius $2$.}

\item \ans{$1$}
\method{$x^2+y^2=2x$ rearranges to $(x-1)^2+y^2=1$, radius $1$.}
\inv{Answering $\sqrt{2}$ (reading the constant $2$ as $r^2$ before completing the square).}

\item \ans{false}
\method{The point $\left(\sqrt2,\sqrt2\right)$ lies on $x^2+y^2=4$ yet has $x\ne\pm2$ and $y\ne\pm2$; the implication fails in both coordinates.}
\inv{Mistaking the circle's bounding box for its points; only the four axis points hit $\pm2$.}

\item \ans{$x=3$}
\method{$A$ and $B$ are at equal height $y=3$, so the perpendicular bisector is the vertical line $x=\frac{1+5}{2}=3$.}
\inv{Writing the horizontal line $y=3$ through both points.}

\end{enumerate}

\section*{Section B \quad Manipulation Drills}

\begin{enumerate}

\item \ans{A) $4$}
\method{On the $x$-axis, $|x|=2|x-6|$ gives the endpoints $x=4$ and $x=12$; the circle between them has centre $(8,0)$ and radius $4$.}
\inv{Quoting the centre distance $8$ (option B) or the gap $6$ (option D).}

\item \ans{B) $y=2$}
\method{$x^2+y^2=x^2+(y-4)^2$ gives $y=2$, the horizontal bisector.}
\inv{Answering $x=2$ (wrong orientation) or $y=4$.}

\item \ans{B) $4\pi$}
\method{The disk has radius $2$, so its area is $\pi r^2=4\pi$.}
\inv{Using $\pi r=2\pi$ (option A) or $r^2=16$ (option D).}

\item \ans{$2\pi$}
\method{Circumference $=2\pi r=2\pi\cdot 1=2\pi$.}
\inv{Mixing circumference with area ($\pi$) or doubling the radius to $4\pi$.}

\item \ans{C) $x^2+y^2<2$}
\method{If $x^2+y^2<1$ then certainly $x^2+y^2<2$; option A is not necessary (a point with $x^2+y^2=0.75$ is inside but not $<\frac12$).}
\inv{Picking a circle whose radius is smaller than the disk (option A) as if it were implied.}

\item \ans{A) $3$}
\method{On the $x$-axis $|x|=3|x-2|$ gives internal $x=\frac32$ and external $x=3$; the right endpoint is $3$.}
\inv{Choosing the left endpoint $\frac32$ or the $y$-absent centre guess.}

\item \ans{B) $x=-\dfrac12$ and $y=\dfrac{\sqrt3}{2}$}
\method{This point satisfies $x^2+y^2=\frac14+\frac34=1$ exactly, so it pins $P$ on the circle. Option A puts $P$ strictly {\em inside}; C and D do not force $x^2+y^2=1$.}
\inv{Choosing a strictly-inside condition (option A) because it is a {\em necessary} symptom of being on the circle.}

\item \ans{$3$}
\method{Distance from $(3,0)$ to the origin is $\sqrt{9}=3$.}
\inv{Answering $6$ (the segment $AB$) or $9$ (the squared distance).}

\item \ans{$(2,3)$}
\method{In $x^2+y^2+2gx+2fy+c=0$, the centre is $(-g,-f)$; with $2g=-4$ and $2f=-6$ we get $(2,3)$.}
\inv{Halving the wrong sign: $(4,6)$ or negating incorrectly.}

\item \ans{$2$}
\method{Completing the square: $(x-2)^2+(y+3)^2=(-9)+4+9=4$, radius $2$.}
\inv{Quoting the constant $9$ as the radius or $r^2$ before shifting.}

\end{enumerate}

\section*{Section C \quad Substitution \& Structure}

\begin{enumerate}

\item \ans{A) $3$}
\method{The distance from the line $x=3$ to $B=(6,0)$ is $|6-3|=3$.}
\inv{Answering $6$ (the gap $A$--$B$, option B).}

\item \ans{A) $(8,0)$}
\method{Endpoints of the locus on the $x$-axis were $x=4$ and $x=12$, so the centre is their midpoint $(8,0)$.}
\inv{Choosing the $B$-side endpoint $(12,0)$ or the midpoint of $AB$, $(3,0)$.}

\item \ans{A) $(2,2)$}
\method{$(2,2)$ has $x^2+y^2=8<9$; the others give $18$, $13$, $16$.}
\inv{Picking a boundary or an exterior point.}

\item \ans{A) $x=3$}
\method{Equal ratios are the perpendicular bisector of $AB$, the line $x=3$.}
\inv{Answering the line through $B$ (option C) or a diagonal.}

\item \ans{B) $4$}
\method{Each of $(1,2)$, $(2,1)$, $(\sqrt5,0)$, $(1,-2)$ satisfies $x^2+y^2=5$.}
\inv{Counting only the obvious axis pair or miscounting by sign.}

\item \ans{B) $8$}
\method{The major axis equals the constant sum $PA+PB=8$.}
\inv{Quoting the focal distance $6$ (option C) or the semi-major axis $4$ (option A).}

\item \ans{A) $a^2+b^2=1$}
\method{Being on the unit circle is exactly $a^2+b^2=1$: necessary (only these points qualify) and sufficient (every such point is on it).}
\inv{Choosing a merely-necessary condition (option B) or a merely-sufficient one.}

\item \ans{C) $0$}
\method{$(x-2)^2+y^2=0$, a single point at $(2,0)$: ``radius'' $0$.}
\inv{Quoting $\sqrt{4}=2$ (option B) without completing the square.}

\end{enumerate}

\section*{Section D \quad Challenge Ramp}

\begin{enumerate}

\item \ans{A) $8$}
\method{Endpoints at $x=4$ and $x=12$ (from $|x|=2|x-6|$), so the chord between them on the $x$-axis is $12-4=8$.}
\inv{Quoting the radius $4$ (option B) or the centre-to-$B$ distance.}

\item \ans{A) $1$}
\method{$(x-2)^2+(y-3)^2=12-4-9+$shift$=1$, radius $1$.}
\inv{Failing to subtract the completed-square constant; $\sqrt{12}$ (option D) is a snare.}

\item \ans{A) $x^2+y^2<\dfrac12$}
\method{This forces $x^2+y^2<1$ strictly, i.e.\ strictly inside the unit disk. Option C allows points outside the unit disk (e.g.\ $x^2+y^2=1.5$).}
\inv{Choosing the weaker disk $x^2+y^2<2$ (which is not sufficient).}

\item \ans{A)}
\method{Expanding $x^2+y^2=9\big((x-2)^2+y^2\big)$ yields $\left(x-\frac94\right)^2+y^2=\frac{9}{16}$, a circle centred at $\left(\frac94,0\right)$ of radius $\frac34$ --- not $(x-2)^2+y^2=9$. Fil's circle is centred at $B$ and ignores $A$ entirely.}
\inv{Believing $(x-2)^2+y^2=9$ (option D) or claiming the ratio locus is a line (option C).}

\item \ans{$x^2+y^2=5$}
\method{Through $(1,2)$ with centre $O$: $r^2=1^2+2^2=5$.}
\inv{Writing $(x-1)^2+(y-2)^2=5$ (putting the given point as the centre).}

\end{enumerate}

\section*{Top 5 Patterns}

\begin{itemize}
\item \textbf{A fixed-distance locus is a circle; square it to find the equation.} Distance $=r$ means $x^2+y^2=r^2$; an offset centre shifts the squares. \seealso{A1, A6, A8, B3, B4, D5}
\item \textbf{$\frac{PA}{PB}=k$ is an Apollonius circle when $k\ne1$.} Square the ratio, collect $x,y$, complete the square; endpoints on the $x$-axis come from $|x|=k|x-b|$. \seealso{A3, B1, B6, C2, D1, D4}
\item \textbf{Equal ratios ($k=1$) are the perpendicular bisector.} The line $\frac{PA}{PB}=1$ is the mirror line of the segment. \seealso{A2, A10, B2, B8, C1, C4}
\item \textbf{Distinguish {\em necessary} from {\em sufficient}.} A condition must be true for the claim (necessary) versus enough to force the claim (sufficient). \seealso{B5, B7, C7, D3}
\item \textbf{A line and a circle can both be ``3 units away''; name both.} Absolute-value loci split into symmetric branches. \seealso{A4, A9}
\end{itemize}

\section*{Common Traps}

\begin{itemize}
\item \textbf{Reading a distance-equality as a point instead of a family.} $PA=PB$ names a whole line, not just the midpoint. \seealso{A2, A10, C4}
\item \textbf{Quoting the radius from $r^2$ before completing the square.} The constant in $x^2+y^2+2gx+2fy+c=0$ is not $r^2$; the centred shift supplies it. \seealso{B10, C8, D2, A8}
\item \textbf{Forgetting the second solution when a modulus meets.} $|x|=k|x-b|$ has an internal and an external endpoint; a flagged ratio answer usually picks the wrong one. \seealso{A3, B6, D1}
\item \textbf{Confusing necessary with sufficient.} ``Must be inside'' and ``is forced on the circle'' are different strengths; one extra decimal of the radius flips the verdict. \seealso{B5, B7, C7, D3}
\item \textbf{Believing a claimed Apollonius equation without testing it.} Expand the squared ratio yourself; a circle centred at one focus is the classic decoy. \seealso{D4}
\end{itemize}

\SpeedClosing{``A locus is a sentence the entire family agrees on; a geometric claim is only as good as the condition you actually invoked.''}

\end{document}